% ============================================================
%  Monte Carlo Engine — Implementation Report
%  Compile with: pdflatex mc_implementation.tex
%  Requires: amsmath, listings, booktabs, geometry, xcolor
% ============================================================
\documentclass[11pt, a4paper]{article}

\usepackage[utf8]{inputenc}
\usepackage[T1]{fontenc}
\usepackage[margin=2.5cm]{geometry}
\usepackage{amsmath, amssymb}
\usepackage{booktabs}
\usepackage{xcolor}
\usepackage{listings}
\usepackage{hyperref}
\usepackage{parskip}
\usepackage{graphicx}
% ── Fortran listings style ───────────────────────────────────
\definecolor{codebg}{RGB}{248,248,248}
\definecolor{codecomment}{RGB}{100,100,100}
\definecolor{codekw}{RGB}{0,80,160}
\definecolor{codestr}{RGB}{160,0,0}

\lstdefinelanguage{Fortran90}{
  morekeywords={program,end,use,implicit,none,integer,double,precision,
                logical,character,allocatable,parameter,real,if,then,else,
                do,while,call,subroutine,function,return,exit,allocate,
                deallocate,write,read,open,close,len,size,max,min,abs,
                sum,sqrt,mod,int,dble,any,trim},
  sensitive=false,
  morecomment=[l]{!},
  morestring=[b]",
  morestring=[b]'
}

\lstset{
  language=Fortran90,
  basicstyle=\ttfamily\small,
  backgroundcolor=\color{codebg},
  commentstyle=\color{codecomment}\itshape,
  keywordstyle=\color{codekw}\bfseries,
  stringstyle=\color{codestr},
  frame=single,
  framesep=4pt,
  rulecolor=\color{gray!40},
  breaklines=true,
  numbers=none,
  captionpos=b,
  xleftmargin=6pt,
  xrightmargin=6pt,
}

% ─────────────────────────────────────────────────────────────
\begin{document}

\title{\textbf{Monte Carlo Simulation of Polyethylene Chains}\\
       \large Implementation Notes}
\author{ManelDC55}
\date{}
\maketitle
\tableofcontents
\newpage

% ============================================================
\section{Serial Monte Carlo Engine}
\label{sec:serial}
% ============================================================

\subsection{The Pivot Move and Rodrigues' Rotation}

The core conformational sampling mechanism is the pivot move, implemented
within the \texttt{rotate\_dihedral} subroutine. Rather than performing local
perturbations, the pivot move selects a random bond along the chain and rotates
the entire subsequent segment (the ``tail'') as a rigid body..

To perform the rotation without distorting bond lengths or bond angles,
Rodrigues' rotation formula is used. Given a vector $\vec{v}$ representing
an atom's position relative to the pivot point, a unit rotation axis $\hat{k}$
(the chosen C--C bond direction), and a rotation angle $\phi$, the rotated vector
is:

\begin{equation}
  \vec{v}_{\,\text{rot}} = \vec{v}\cos\phi
    + (\hat{k} \times \vec{v})\sin\phi
    + \hat{k}\,(\hat{k} \cdot \vec{v})(1 - \cos\phi)
  \label{eq:rodrigues}
\end{equation}

The formula decomposes the rotation into three contributions: a scaled version
of the original vector, a cross-product term that generates the perpendicular
component, and a projection term that preserves the component parallel to the
axis.

When explicit hydrogens are included (all-atom mode), the rotation must also
displace the two hydrogens attached to each rotating carbon. Because of the
$sp^3$ hybridisation, both hydrogens must rotate with exactly the same
axis and angle as their parent carbon; any deviation would alter C--H bond
lengths or C--C--H angles, producing non-physical geometries that would be
immediately rejected by the Metropolis criterion.

The implementation identifies each hydrogen by checking whether the distance
to a rotating carbon falls below the C--H bond threshold of $1.2\,\text{\AA}$,
then applies Eq.~\eqref{eq:rodrigues} with the same pivot, axis, and angle:

\begin{lstlisting}[caption={Pivot move --- hydrogen rotation in AA mode.}]
if (explicit_h) then
  ! Hydrogens are stored after the carbon backbone in the coords array
  do i = n_carbons + 1, n_atoms
    do j = k + 1, n_carbons
      ! Identify if hydrogen i is bonded to the moving carbon j
      v = coords(i, :) - coords(j, :)
      if (vnorm(v) < 1.2d0) then
        ! Apply Rodrigues' rotation relative to the pivot point
        v = coords(i, :) - pivot
        dot_uv = sum(axis * v)
        v_rot = v * cos_p + cross(axis, v) * sin_p &
              + axis * dot_uv * (1.0d0 - cos_p)
        coords_new(i, :) = pivot + v_rot
        exit  ! Hydrogen found, move to the next atom
      end if
    end do
  end do
end if
\end{lstlisting}

% ------------------------------------------------------------
\subsection{Convergence Detection: the Geweke $z$-Test}
\label{sec:geweke}
% ------------------------------------------------------------

Detecting equilibration of a Markov chain is non-trivial because successive
samples are highly correlated. The Geweke diagnostic was used, as it is designed specifically for Markov chain
output and accounts for this autocorrelation.

The test compares the mean energy in the first $f_A = 10\%$ of an accumulated
buffer of $N = 300$ energy samples against the mean in the last $f_B = 50\%$:

\begin{equation}
  z = \frac{\bar{x}_A - \bar{x}_B}{\sqrt{SE_A^2 + SE_B^2}}
  \label{eq:geweke}
\end{equation}

Here, $SE_A$ and $SE_B$ represent the Standard Errors of the mean energy
for window A and window B, respectively. They measure the statistical uncertainty 
of our calculated averages. 

In a dataset of independent measurements, calculating the standard error is 
straightforward. However, Monte Carlo trajectories are autocorrelated---each 
new geometry is simply the previous one with a slight rotation. If we computed 
$SE_A$ and $SE_B$ directly from the raw sample variance, the autocorrelation 
would cause us to drastically underestimate the error, tricking the algorithm 
into declaring a false equilibrium.

To solve this, the code computes $SE_A$ and $SE_B$ using batch means. 
Instead of treating all $n$ samples in a window individually, the window is 
divided into $\lfloor\sqrt{n}\rfloor$ discrete batches or blocks. By calculating 
the variance across the averages of these batches, the local autocorrelation is 
"smoothed out," providing a consistent and conservative estimator of the true 
long-run variance.

Under the null hypothesis of stationarity, $z$ follows a standard normal
distribution. Equilibrium is accepted when $|z| < z_{\text{crit}} = 1.96$
(95\% confidence level) on $n_{\text{consec}} = 3$ consecutive evaluations,
preventing false positives caused by temporary plateaus. The test is re-evaluated
every \texttt{eval\_freq} synchronisation points once the buffer is full.

\begin{lstlisting}[caption={Batch-means SE and Geweke $z$-statistic.}]
! Batch means SE for window A (first fA% of buffer)
nA    = max(2, n_geweke * fA_pct / 100)   ! = 30 samples
bA    = max(2, int(sqrt(dble(nA))))        ! number of batches
bsA   = nA / bA                            ! batch size
seA_E = 0.0d0
do ib = 1, bA
  bm    = sum(tmp_E((ib-1)*bsA+1 : ib*bsA)) / dble(bsA)
  seA_E = seA_E + (bm - meanA_E)**2
end do
seA_E = seA_E / dble(bA * (bA - 1))       ! SE^2 of window A mean

! Geweke z-statistic
z_E = abs(meanA_E - meanB_E) / sqrt(seA_E + seB_E)

if (z_E < z_crit) then   ! z_crit = 1.96 (95% confidence)
  consec_passes(task) = consec_passes(task) + 1
  if (consec_passes(task) >= n_consec) equil_done(task) = .true.
else
  consec_passes(task) = 0  ! reset counter on failure
end if
\end{lstlisting}

The radius of gyration $R_g$ is monitored with its own Geweke $z$-statistic and
reported in the output, but is deliberately excluded from the stopping criterion.
In practice, $R_g$ was found to equilibrate at a different rate than the energy
and would prevent convergence from being declared even when the energy had
clearly stabilised.

% ============================================================
\newpage
\section{Parallel Equilibration via MPI}
\label{sec:mpi}
% ============================================================

\subsection{Strategy: Collaborative Population Equilibration}

Rather than assigning a single worker to equilibrate each configuration type
independently, a group of \texttt{workers\_per\_equil} MPI workers is assigned
to each configuration type simultaneously. Each worker starts from the same
initial geometry but with a different random seed
($\texttt{rng\_seed} + \texttt{rank} \times 104729$), ensuring that trajectories
immediately diverge and explore distinct regions of conformational space.

Every \texttt{sync\_interval} MC steps, a synchronisation protocol is executed:
each worker sends its current energy and $R_g$ to the master, which selects the
most representative configuration via a Boltzmann-weighted roulette wheel and
broadcasts those coordinates back to the non-winning workers in the group. The
winning worker continues its trajectory uninterrupted, preserving the best
conformational path found so far.

\subsection{Replica Selection: Boltzmann Roulette Wheel}

At each synchronisation point, the master assigns a statistical weight to each
replica $i$ based on its total energy $E_i$:

\begin{equation}
  w_i = \exp\!\left(-\frac{E_i - E_{\min}}{k_B \, T_{\text{virt}}}\right)
  \label{eq:boltzmann_weight}
\end{equation}

where $E_{\min}$ is the lowest energy in the current group and $T_{\text{virt}}$
is a virtual temperature — a purely algorithmic parameter with no
physical meaning. It controls how aggressively low-energy replicas are favoured:

\begin{itemize}
  \item $T_{\text{virt}} \to 0$: deterministic selection of the minimum-energy replica.
  \item $T_{\text{virt}} \to \infty$: uniform selection, ignoring energies.
\end{itemize}

A uniform random number then samples this distribution via roulette-wheel
selection, and the selected replica's coordinates are sent to the non-winning
workers:

\begin{lstlisting}[caption={Boltzmann roulette wheel selection.}]
E_min_grp = minval(grp_energies)
do w = 1, workers_per_equil
  grp_boltz(w) = exp(-(grp_energies(w) - E_min_grp) / (kb * T_virt))
end do
tw = sum(grp_boltz)
call random_number(r_pick)
accum = 0.0d0
best_local = workers_per_equil
do w = 1, workers_per_equil
  accum = accum + grp_boltz(w) / tw
  if (r_pick <= accum) then
    best_local = w
    exit
  end if
end do
\end{lstlisting}

\subsection{MPI Communication Protocol}

The synchronisation uses only point-to-point MPI operations (\texttt{MPI\_Send}
/ \texttt{MPI\_Recv}) to avoid the need for temporary communicators or collective
operations across worker subgroups. At each sync point, the master loops over
all $K$ workers in the active equilibration group following the protocol
described in Table~\ref{tab:mpi_protocol}.

\begin{table}[ht]
\centering
\caption{Point-to-point MPI synchronisation protocol at each \texttt{sync\_interval}.}
\label{tab:mpi_protocol}
\begin{tabular}{clll}
\toprule
Step & Direction & Tag & Content \\
\midrule
1 & worker $\to$ master & \texttt{TAG\_SYNC\_OBS}    & $[E_{\text{total}},\, R_g^2]$ \\
2 & master $\to$ worker & \texttt{TAG\_SYNC\_CTRL}   & $[\texttt{ctrl\_signal},\, \texttt{best\_local\_idx}]$ \\
3 & worker $\to$ master & \texttt{TAG\_EQUIL\_COORDS}& full coordinate array \\
4 & master $\to$ non-winners & \texttt{TAG\_SYNC\_COORDS} & winning coordinates \\
\bottomrule
\end{tabular}
\end{table}

Once the Geweke test (Section~\ref{sec:geweke}) declares equilibrium, the master
saves the equilibrated geometry to \texttt{../results/equilibrated\_cX.xyz} and
unlocks the 10 production jobs for that configuration type in the task queue.


% ============================================================

% ============================================================
\newpage
\section{Parallelisation Efficiency and Sampling Results}
\label{sec:efficiency}
% ============================================================

\subsection{Metrics and Methodology}

This section evaluates the performance of the collaborative parallel implementation compared to the serial baseline. The system studied the configuration 1 chain of 500 carbons at $T = 300\,\text{K}$ including hydrogens 

We evaluate efficiency based on the time to reach specific energy thresholds ($E_{total} < 110, 100, \text{ and } 90$ kcal/mol). The speedup is defined as $S(w) = \text{Steps}_{\text{serial}} / \text{Steps}_{\text{parallel}}(w)$.

\subsection{Convergence Efficiency and Superlinear Speedup}

Table~\ref{tab:energy_thresholds} summarises the number of MC steps required by different worker configurations to reach the target energy levels.

\begin{table}[ht]
\centering
\caption{Sampling efficiency: MC steps and speedup to reach energy thresholds.}
\label{tab:energy_thresholds}
\begin{tabular}{lcrrc}
\toprule
Threshold & Workers ($w$) & Steps Required & Speedup $S(w)$ & Efficiency $S(w)/w$ \\
\midrule
$E < 110$ kcal/mol & 1 & 590,000 & 1.00 & 1.00 \\
                   & 2 & 70,000  & 8.43 & 4.21 \\
                   & 4 & 140,000 & 4.21 & 1.05 \\
                   & 8 & 50,000  & 11.80 & 1.47 \\
\midrule
$E < 100$ kcal/mol & 1 & 1,180,000 & 1.00 & 1.00 \\
                   & 2 & 350,000   & 3.37 & 1.68 \\
                   & 4 & 210,000   & 5.62 & 1.40 \\
                   & 8 & 50,000    & 23.60 & 2.95 \\
\midrule
$E < 90$ kcal/mol  & 1 & 2,400,000 & 1.00 & 1.00 \\
                   & 2 & 1,780,000 & 1.35 & 0.67 \\
                   & 4 & 570,000   & 4.21 & 1.05 \\
                   & 8 & 160,000   & 15.00 & 1.87 \\
\bottomrule
\end{tabular}
\end{table}

\begin{figure}[ht]
  \centering
  \includegraphics[width=\textwidth]{convergence_speedup.png}
  \caption{Convergence speedup analysis. \textit{Left}: evolution of $E_\text{total}$ for each worker configuration. \textit{Centre}: MC steps required to reach each energy threshold. \textit{Right}: speedup $S(w)$ relative to the serial run; the dashed line shows
  ideal linear scaling.}
  \label{fig:convergence_speedup}
\end{figure}


\subsection{Discussion}

The results demonstrate a clear superlinear speedup ($S(w) > w$) in almost all parallel configurations. This phenomenon is a direct consequence of the collaborative sampling strategy implemented via the Boltzmann roulette wheel.

In the serial run, the polymer chain often becomes trapped in metastable states (local energy minima). The parallel implementation, by maintaining a population of $w$ independent trajectories, significantly increases the probability that at least one replica discovers a more favorable region of the phase space. Once a lower-energy conformation is found, the synchronization protocol "rescues" the rest of the population, broadcasting the winning coordinates to all workers. This creates a collective escape mechanism that reduces the number of steps required to converge by a factor far greater than the number of processors.

For instance, at the 100 kcal/mol threshold, the 8-worker group reaches the target in only 50,000 steps, compared to the 1.18 million steps required by the serial run, resulting in a speedup of \textbf{23.60x}. 

We also observe stochastic fluctuations typical of Monte Carlo methods, such as the anomalous efficiency at $w=2$ for the 110 kcal/mol threshold. As the energy targets become more stringent (90 kcal/mol), the speedup values tend to stabilize, as finding further conformational improvements becomes more difficult even for a larger population. 

\end{document}